\section{Desain}\label{sec:Desain}

\subsection{Tech Stack}\label{sub:Tech Stack}

Aplikasi dirancang menggunakan Go untuk bagian parser, model puzzle,
dan solver. Untuk bonus GUI, aplikasi menggunakan backend HTTP Go dan
frontend React berbasis Vite.

\begin{enumerate}
    \item Bahasa pemrograman solver: Go 1.26.2.
    \item Backend GUI: Go HTTP server.
    \item Frontend GUI: React dan Vite.
    \item Kontainerisasi: Docker multi-stage build.
\end{enumerate}

\subsection{Struktur Program}\label{sub:Struktur Program}

Struktur utama repository adalah sebagai berikut.
\begin{enumerate}
    \item \texttt{src/main.go}: entry point CLI.
    \item \texttt{src/parser/}: parser dan validasi berkas input.
    \item \texttt{src/game/}: representasi state permainan, aturan
        slide, dan pembangkitan successor.
    \item \texttt{src/solver/}: implementasi UCS, GBFS, A*, BFS, DFS,
        priority queue, dan heuristic.
    \item \texttt{src/searchtrace/}: struktur data trace pencarian
        untuk kebutuhan visualisasi GUI.
    \item \texttt{src/web/backend/}: backend HTTP untuk menerima file
        input dan menjalankan solver.
    \item \texttt{src/web/frontend/}: antarmuka React untuk upload
        testcase, kontrol solver, playback solusi, dan graf pencarian.
    \item \texttt{test/}: tempat penyimpanan testcase.
    \item \texttt{doc/}: spesifikasi dan laporan.
\end{enumerate}

\subsection{Representasi Data}\label{sub:Representasi Data}

Konfigurasi papan direpresentasikan oleh struktur \texttt{BoardConfig}.
Struktur ini menyimpan ukuran papan, grid karakter, matriks biaya,
posisi awal, posisi tujuan, dan posisi checkpoint.

\begin{lstlisting}[language=Go, basicstyle=\ttfamily\small, breaklines=true]
type BoardConfig struct {
    N, M           int
    Grid           [][]rune
    Cost           [][]int
    StartRow       int
    StartCol       int
    GoalRow        int
    GoalCol        int
    NumCheckpoints int
    CheckpointPos  map[int][2]int
}
\end{lstlisting}

Status pencarian direpresentasikan oleh \texttt{game.Node}. Selain
posisi pemain, state juga menyimpan checkpoint berikutnya yang harus
dicapai. Dengan demikian, posisi yang sama dapat menjadi state yang
berbeda jika progres checkpoint-nya berbeda.

\begin{lstlisting}[language=Go, basicstyle=\ttfamily\small, breaklines=true]
type Node struct {
    Row, Col   int
    NextTarget int
}
\end{lstlisting}

\subsection{Model Gerakan}\label{sub:Model Gerakan}

Gerakan pada puzzle tidak berpindah satu petak, melainkan slide
sampai petak berikutnya adalah dinding. Fungsi slide memeriksa arah
gerak, menambahkan biaya petak yang dilewati, memperbarui checkpoint,
dan menolak langkah yang keluar papan, melewati lava, atau melewati
checkpoint yang belum boleh dilewati.

Setiap successor yang valid direpresentasikan sebagai \texttt{Edge}.
\begin{lstlisting}[language=Go, basicstyle=\ttfamily\small, breaklines=true]
type Edge struct {
    State     Node
    Cost      int
    Direction Dir
    Path      []Pos
}
\end{lstlisting}

\subsection{Desain Solver}\label{sub:Desain Solver}

Solver menggunakan satu kerangka pencarian umum, yaitu
\texttt{graphSearch}. Perbedaan UCS, GBFS, dan A* hanya terletak pada
cara menghitung prioritas simpul.

\begin{table}[H]
    \centering
    \begin{tabular}{|c|c|}
        \hline
        \textbf{Algoritma} & \textbf{Prioritas} \\
        \hline
        UCS & $g(n)$ \\
        \hline
        GBFS & $h(n)$ \\
        \hline
        A* & $g(n) + h(n)$ \\
        \hline
    \end{tabular}
    \caption{Fungsi prioritas pada solver}
\end{table}

Frontier disimpan dalam priority queue. Program juga menyimpan
\texttt{bestCost} untuk setiap state agar state yang pernah dicapai
dengan biaya lebih murah tidak diproses ulang menggunakan biaya yang
lebih mahal.

\subsection{Desain BFS dan DFS}\label{sub:Desain BFS DFS}

BFS dan DFS diimplementasikan dengan kerangka pencarian tidak berbobot.
Keduanya tetap memakai fungsi \texttt{GetSuccessors} yang sama dengan
UCS, GBFS, dan A*. Perbedaannya terletak pada struktur frontier.

\begin{table}[H]
    \centering
    \begin{tabular}{|c|c|L{7cm}|}
        \hline
        \textbf{Algoritma} & \textbf{Frontier} & \textbf{Karakteristik} \\
        \hline
        BFS & Queue & Menemukan solusi dengan jumlah slide minimum, tetapi tidak menjamin total cost minimum. \\
        \hline
        DFS & Stack & Menelusuri cabang sedalam mungkin dan tidak menjamin optimalitas. \\
        \hline
    \end{tabular}
    \caption{Desain frontier BFS dan DFS}
\end{table}

Pada BFS dan DFS, state yang sudah pernah dikunjungi tidak dimasukkan
kembali ke frontier untuk mencegah siklus. Total cost tetap dihitung
sepanjang lintasan yang ditemukan agar hasilnya dapat dibandingkan
dengan algoritma berbobot.

\subsection{Cara Kompilasi dan Menjalankan Program}\label{sec:Cara Kompilasi dan Menjalankan Program}

Program membutuhkan Go untuk menjalankan CLI dan backend. Frontend
GUI membutuhkan Node.js dan npm jika dijalankan tanpa Docker. Semua
perintah berikut dijalankan dari root repository.

\noindent Untuk melakukan kompilasi executable:
\begin{lstlisting}[language=bash, basicstyle=\ttfamily\small, breaklines=true]
go build -o bin/ice-sliding-solver ./src
\end{lstlisting}

\noindent Untuk menjalankan hasil kompilasi:
\begin{lstlisting}[language=bash, basicstyle=\ttfamily\small, breaklines=true]
./bin/ice-sliding-solver
\end{lstlisting}

\noindent Untuk menjalankan CLI:
\begin{lstlisting}[language=bash, basicstyle=\ttfamily\small, breaklines=true]
go run ./src
\end{lstlisting}

Setelah solver selesai, CLI menawarkan penyimpanan hasil ke berkas
\texttt{.txt}. Berkas hasil berisi algoritma, heuristic jika ada,
status solusi, urutan gerakan, total cost, jumlah iterasi, waktu
eksekusi, dan visualisasi papan per langkah.

\noindent Untuk menjalankan backend GUI tanpa Docker:
\begin{lstlisting}[language=bash, basicstyle=\ttfamily\small, breaklines=true]
go run ./src/web/backend
\end{lstlisting}

\noindent Untuk menjalankan frontend GUI tanpa Docker:
\begin{lstlisting}[language=bash, basicstyle=\ttfamily\small, breaklines=true]
cd src/web/frontend
npm install
npm run dev
\end{lstlisting}

\noindent Untuk menjalankan GUI menggunakan Docker:
\begin{lstlisting}[language=bash, basicstyle=\ttfamily\small, breaklines=true]
docker compose down
docker compose up --build
\end{lstlisting}

\noindent GUI dapat diakses melalui alamat berikut.
\begin{lstlisting}[language=bash, basicstyle=\ttfamily\small, breaklines=true]
http://localhost:8080
\end{lstlisting}
